% Generated draft: requires mathematical and source review.
\documentclass[11pt]{article}
\usepackage[margin=25mm]{geometry}
\usepackage{fontspec}
\setmainfont{Latin Modern Roman}
\usepackage{amsmath,amssymb,mathtools,unicode-math}
\setmathfont{Latin Modern Math}
\usepackage{enumitem,needspace,xurl,hyperref}
\setlength{\parindent}{0pt}
\setlength{\parskip}{5pt}
\setlength{\emergencystretch}{3em}
\newcommand{\N}{\mathbb{N}}
\newcommand{\Z}{\mathbb{Z}}
\newcommand{\Q}{\mathbb{Q}}
\newcommand{\R}{\mathbb{R}}
\newcommand{\C}{\mathbb{C}}
\newcommand{\F}{\mathbb{F}}
\newcommand{\A}{\mathbb{A}}
\newcommand{\PP}{\mathbb{P}}
\newcommand{\E}{\mathbb{E}}
\newcommand{\eps}{\varepsilon}
\newcommand{\dd}{\,\mathrm{d}}
\begin{document}
\section*{UnsolvedMath ID: 30004350}
\textbf{Source problem code:} catalogue entry
\subsection{Definitions and mathematical statement}
The supplied record concerns a reductive group \(G\) over \(\Q\), a finite set \(S\) of primes, and an irreducible cuspidal automorphic representation \(\pi\) having nonzero \(K\)-invariants, written \(\pi^K\ne\{0\}\), and satisfying the supplied condition that \(\pi\) is not \(S\)-CAP. The subgroup \(K\), the definition of \(S\)-CAP, the meaning of an \(S\)-multiplier, and the admissible test functions are not specified in the supplied record.

The asserted spectral-isolation property can therefore be transcribed only schematically. Write \(\mathcal H\) for the unspecified ambient automorphic Hilbert space, \(\mathcal H_{\mathrm{cusp}}\) for its cuspidal subspace, and \(V_\pi\subseteq\mathcal H_{\mathrm{cusp}}\) for a realization of \(\pi\). Here \(R(f)\) denotes the convolution operator associated with a test function \(f\), and \(\mu\star f\) denotes the multiplier action written as convolution in the source. These symbols describe the roles required by the supplied claim; they do not resolve the missing analytic definitions.

With those qualifications, the supplied assertion is that for every such \(G,S,\pi\), there exists an \(S\)-multiplier \(\mu_\pi\) such that, for every admissible test function \(f\),
\[
\left.R(\mu_\pi\star f)\right|_{\mathcal H_{\mathrm{cusp}}^{\perp}}=0,
\qquad
\left.R(\mu_\pi\star f)\right|_{V_\pi}
=
\left.R(f)\right|_{V_\pi}.
\]
The orthogonal complement is taken inside \(\mathcal H\). The existential quantifier for \(\mu_\pi\) precedes the universal quantifier for \(f\), so the multiplier is required to work simultaneously for all admissible test functions. The record does not determine whether the second identity is intended on a chosen copy of \(\pi\) or its entire isotypic component.

This is a qualified transcription, not a complete theorem statement. In particular, the ambient automorphic quotient, central-character or central-quotient conventions, operator domains, and test-function and multiplier spaces require specification. The supplied registry labels the item open and its statement status unrecoverable. Its separate statement-verification note instead identifies the extraction as Theorem 1 of the cited report, a proved spectral-isolation result rather than an open problem. The earlier supplied literature assessment describes a truncated fragment beginning with ``Conjecture'', which does not match the displayed assertion. These conflicting descriptions remain unresolved here.
\subsection{Short English statement}
The supplied assertion asks for a multiplier that preserves the action of every admissible test function on a specified non-\(S\)-CAP cuspidal representation while making that action vanish on the orthogonal complement of the cuspidal spectrum. Essential definitions and analytic conventions are missing. The supplied open label conflicts with a verification note identifying the result as a proved theorem; neither description is independently verified here.
\subsection{Sources}
\begin{itemize}
\item Supplied local UnsolvedMath registry record, permanent ID \(30004350\), title ``Cuspidal Spectral-Multiplier Theorem'', exact source problem code \texttt{catalogue entry}, in the Oberwolfach Reports Open Problems collection. No individual proposer is supplied; the proposed year is recorded as 2019.
\item Supplied report citation: \emph{Modular Forms (2020)}, Oberwolfach Reports, DOI \texttt{10.4171/owr/2019/57}, pp.\ 3529--3606. Supplied source URL and literature-source link, labeled ``original OWR report'': \url{https://doi.org/10.4171/owr/2019/57}. The report title year and DOI year are retained as supplied.
\item The literature assessment was checked on \texttt{2026-08-22} according to the supplied record. This is supplied dated curation, not an independently verified literature assessment. It classifies the item as open under a rule for unclear records because it describes the extraction as a truncated fragment beginning with ``Conjecture''.
\item The supplied statement-verification note identifies the extraction as Theorem 1 in the official report and describes it as a proved spectral-isolation result for non-\(S\)-CAP cuspidal representations. It says that no separate clean statement was promoted and that the original record was retained for provenance. This source-identification claim has not been independently verified here.
\end{itemize}
\end{document}
